% =============================================================================
%  Final Lab - Demonstration and Extension
%  Plus the material that closes the whole report.
% =============================================================================

\chapter{Final Lab --- Demonstration and Extension}

\section{Executive summary}
% ONE PAGE. What did you build, and what does it do?
%
% Write this for someone who has not read the other twelve chapters - a future
% employer, or yourself in two years. What is the machine, what instruction set
% does it implement, how is it organised, what program does it run, and what
% does that program compute?
%
% Write it LAST, when everything else is finished.

\section{Resource and timing summary}

% One table, every component you built. Pull the numbers from the E4 and E6
% sections of your earlier chapters. This is the first time they sit together,
% and the comparison is the point.

\begin{table}[htbp]
\centering
\caption{Resources and timing across the whole design.}
\label{tab:final-resources}
\begin{tabular}{lrrrr}
\toprule
component & ALMs & registers & block memory bits & Fmax (MHz) \\
\midrule
\texttt{TransLatch} (Lab 1) & & & & \\
\texttt{RegN} (Lab 2) & & & & \\
\texttt{Mux3x32} (Lab 3) & & & & \\
\texttt{ImmedXtnd} (Lab 4) & & & & \\
\texttt{ALU} (Lab 5) & & & & \\
\texttt{RegFile} (Lab 6) & & & & \\
\texttt{NIAG} (Lab 8) & & & & \\
\texttt{ROM} (Lab 9) & & & & \\
\texttt{Decoder} (Lab 11) & & & & \\
\midrule
\texttt{CPU} (Lab 12) & & & & \\
\bottomrule
\end{tabular}
\end{table}

% Two sentences, required:
%   (a) Which single component dominates the area, and why?
%   (b) Compare the CPU's Fmax with the ALU's from Lab 5. Did adding everything
%       else make the machine slower, and if not, why not?

\section{Extension}
% State which track you chose, then complete the matching subsection and DELETE
% the other two.

\subsection*{Track A --- a new instruction}
% (a) Which instruction, and its encoding.
% (b) Your new control-table row, as a table, showing only the differences
%     from CTRL_DEFAULT.
% (c) Why that extension choice - sign or zero - and what the alternative would
%     have done to the test program.
% (d) The verification: the decoder-level check AND the CPU-level run. Explain
%     why either one alone would be insufficient.
% (e) IF YOU DID SUB: what does your CPU's SUB actually compute, and why is it
%     not what real MIPS computes? What in the datapath fixed that, and which
%     earlier instructions hid the problem from you?

\subsection*{Track B --- performance measurement}
% (a) The measured instruction mix, as a table.
% (b) The identity check: the four classes must sum to the cycle count. State
%     yours, and say what a mismatch would have meant.
% (c) CPI as built, and CPI under the two hypothetical machines.
% (d) REQUIRED: which hypothetical hurts more, and why? Answer in terms of the
%     measured mix, not in general terms. This is Amdahl's law applied to a
%     number you measured yourself.
% (e) Name one change to this program - not the hardware - that would change
%     the mix substantially, and predict its effect on both hypotheticals.

\subsection*{Track C --- pipeline and hazards}
% (a) What you built, and the straight-line program that demonstrates stage 2
%     working.
% (b) THE HAZARD ANALYSIS. This is the graded part. For each of:
%         control hazard, data hazard, structural hazard, reset state
%     say whether it is PRESENT in your design, and:
%         - if present: a CONCRETE instruction pair from the Lab 12 program
%           that triggers it, one named real-world fix, and what that fix costs
%         - if absent: why, and what design change would introduce it
% (c) REQUIRED: the control hazard is the reason real MIPS has a branch delay
%     slot. Explain that connection in your own words.
% (d) An honest statement of what does not work, and why.

\section{Reflection --- what I would design differently}
% Not a summary. A judgement.
%
% Pick two or three specific decisions made for you during the semester - the
% eight-instruction subset, big-endian byte order, Asel = rt, the two-phase
% timing, LUI as a plain immediate load, the read-old register file - and say,
% for each:
%     what the decision bought
%     what it cost
%     what you would choose now, and why
%
% "I would add more instructions" is not an answer. Name a decision, name its
% trade-off, take a position.

\section{Discussion}
% One last time: the most valuable thing you learned that you did NOT expect to
% learn. Not the material - the method, the debugging, the testing, or the way
% a specification turns into hardware.

\appendix
\chapter{Code}
% Every .vhd file you wrote, via \vhdlfile. Components you did not write
% yourself do not belong here - say which those are in the Executive Summary.
